\documentclass[conference]{IEEEtran}
\IEEEoverridecommandlockouts

\usepackage[T1]{fontenc}
\usepackage[utf8]{inputenc}
\usepackage{lmodern}

\usepackage{cite}
\usepackage{amsmath,amssymb,amsfonts}
\usepackage{algorithmic}
\usepackage{graphicx}
\usepackage{textcomp}
\usepackage{xcolor}
\usepackage{url}
\usepackage{datetime}
\usepackage{microtype}
\usepackage{parskip}

\newdateformat{schoolyear}{%
  \ifthenelse{\month > 9}
    {\the\year/\the\numexpr\year+1-2000\relax}
    {\the\numexpr\year-1\relax/\the\numexpr\year-2000\relax}
}

\begin{document}

\title{Assignment \#4:\\Improving Face Recognition with Deep Detectors and Embeddings}

\author{
\IEEEauthorblockN{Dino Dzaferagic}
\IEEEauthorblockA{\textit{IBB \schoolyear\today, FRI, UL}\\
dd5102@student.uni-lj.si}
}

\maketitle

% Per instructions: Abstract, Introduction, Related Work omitted.

% ########################################
\section{Methodology}
The goal of Assignment~4 is to improve the face recognition system from Assignment~3 by incorporating deep-learning methods.
The evaluated pipeline consists of: (i) face detection and cropping, (ii) feature extraction, and (iii) gallery--probe matching with cosine similarity and CMC evaluation.
We compare classical descriptors (LBP\_ms, HOG, Dense SIFT) with deep face embeddings (FaceNet and ArcFace), where all embedding vectors are L2-normalized.
Here, LBP\_ms denotes multi-scale LBP, arcface\_insightface denotes ArcFace embeddings from the InsightFace model, and facenet\_vggface2 denotes FaceNet embeddings with VGGFace2-pretrained weights.
CMC summarizes identification accuracy as a function of rank; Rank-1 and Rank-5 report the percentage of probes whose true identity appears within the top-1 or top-5 matches.
Each method is evaluated on (a) whole images and (b) a detection-based pipeline, where faces are cropped using the best detector.

% ########################################
\section{Experiments}
\textbf{Dataset and protocol.} CelebA-HQ-small with 50 identities. The fixed split contains 475 training images and 412 test images.
The test set is used for evaluation with 100 gallery images and 312 probe images (closed-set identification).

\textbf{Detection.} Three detectors are evaluated: Viola--Jones (with grid-search tuning over \texttt{scaleFactor}, \texttt{minNeighbors}, and \texttt{minSize}), YOLO face detector (YOLOv8 face weights), and the InsightFace detector (buffalo\_l).
Detection is measured with mean IoU (against ground-truth face boxes) and detection rate.
The best detector (by mean IoU) is then used to crop faces for the pipeline evaluation.

\textbf{Recognition.} We evaluate three classical feature extractors (LBP\_ms, HOG, Dense SIFT) and two deep embedding models (ArcFace from InsightFace ONNX and FaceNet/VGGFace2).
Similarity is computed with cosine distance and performance is summarized using CMC with Rank-1 and Rank-5.

% ########################################
\section{Results and Discussion}

\subsection{Results}
\vspace{-0.2em}
Table~\ref{tab:detection} shows that YOLO achieves the highest IoU while maintaining a high detection rate, so its crops are used in the pipeline.
Table~\ref{tab:recognition} and Figs.~\ref{fig:cmc_whole_classical}--\ref{fig:cmc_pipeline_facenet} summarize identification accuracy on whole images and cropped faces.
Deep embeddings dominate classical descriptors, with ArcFace and FaceNet reaching near-perfect Rank-1/5, while classical descriptors benefit from tighter crops.

\begin{table}[t]
\caption{Detection performance on the test set. YOLO achieves the highest mean IoU and is therefore used for pipeline crops.}
\centering
\scriptsize
\begin{tabular}{lcc}
\textbf{Detector} & \textbf{Mean IoU} & \textbf{Detection rate [\%]} \\
\hline
Viola--Jones & 0.686 & 96.84 \\
YOLO (deep)  & \textbf{0.841} & 96.12 \\
InsightFace  & 0.435 & 48.06 \\
\end{tabular}
\label{tab:detection}
\vspace{-1.0em}
\end{table}

\vspace{-0.4em}
\begin{table}[t]
\caption{Recognition performance (CMC). Rank-1/5 on whole images and pipeline (YOLO crops). Pipeline uses 310/312 probes due to 2 failed detections.}
\centering
\scriptsize
\begin{tabular}{lcccc}
\textbf{Feature} &
\textbf{Whole R1} & \textbf{Whole R5} &
\textbf{Pipe R1} & \textbf{Pipe R5} \\
\hline
LBP\_ms                 & 11.2 & 30.1 & 15.2 & 36.1 \\
HOG                    & 15.1 & 38.1 & 19.4 & 41.9 \\
Dense SIFT             & 19.9 & 42.3 & 25.2 & 45.5 \\
ArcFace (InsightFace)  & \textbf{100.0} & \textbf{100.0} & 97.1 & 98.4 \\
FaceNet (VGGFace2)     & 95.2 & \textbf{100.0} & \textbf{97.4} & \textbf{98.7} \\
\end{tabular}
\label{tab:recognition}
\vspace{-1.0em}
\end{table}

% --- Combined classical CMCs (whole + pipeline) ---
\begin{figure}[t]
\centering
\includegraphics[width=0.98\columnwidth]{figs/cmc_classical_whole_combined.png}
\vspace{-0.6em}
\caption{CMC curves on whole images for classical features (LBP\_ms, HOG, Dense SIFT).}
\label{fig:cmc_whole_classical}
\vspace{-0.9em}
\end{figure}

\begin{figure}[t]
\centering
\includegraphics[width=0.98\columnwidth]{figs/cmc_classical_pipeline_combined.png}
\vspace{-0.6em}
\caption{CMC curves for the pipeline (YOLO crops) for classical features (LBP\_ms, HOG, Dense SIFT).}
\label{fig:cmc_pipeline_classical}
\vspace{-0.9em}
\end{figure}

% --- Deep models (standalone plots) ---
\begin{figure}[t]
\centering
\includegraphics[width=0.98\columnwidth]{figs/cmc_whole_arcface_insightface.png}
\vspace{-0.6em}
\caption{CMC curve for ArcFace on whole images.}
\label{fig:cmc_whole_arcface}
\vspace{-0.7em}
\end{figure}

\begin{figure}[t]
\centering
\includegraphics[width=0.98\columnwidth]{figs/cmc_pipeline_arcface_insightface.png}
\vspace{-0.6em}
\caption{CMC curve for ArcFace on the YOLO-cropped pipeline.}
\label{fig:cmc_pipeline_arcface}
\vspace{-0.7em}
\end{figure}

\begin{figure}[t]
\centering
\includegraphics[width=0.98\columnwidth]{figs/cmc_whole_facenet_vggface2.png}
\vspace{-0.6em}
\caption{CMC curve for FaceNet (VGGFace2) on whole images.}
\label{fig:cmc_whole_facenet}
\vspace{-0.7em}
\end{figure}

\begin{figure}[t]
\centering
\includegraphics[width=0.98\columnwidth]{figs/cmc_pipeline_facenet_vggface2.png}
\vspace{-0.6em}
\caption{CMC curve for FaceNet (VGGFace2) on the YOLO-cropped pipeline.}
\label{fig:cmc_pipeline_facenet}
\vspace{-0.7em}
\end{figure}

\subsection{Discussion}
YOLO achieves the best localization accuracy (Table~\ref{tab:detection}), so it is used for face cropping in the pipeline.
Classical descriptors benefit from tighter face crops (Table~\ref{tab:recognition}, Fig.~\ref{fig:cmc_pipeline_classical}) because background and pose variability are reduced.
Deep embeddings dominate performance: ArcFace reaches perfect Rank-1 on whole images and remains near-perfect in the pipeline, while FaceNet also achieves very high accuracy.
The InsightFace detector produced a much lower detection rate, likely due to default thresholds or detector configuration mismatch for this dataset.
The slight differences between whole-image and pipeline results are expected due to occasional detection failures (310 probes used) and sensitivity to crop/resizing.

% ########################################
\section{Conclusion}
This assignment improved the baseline face recognition system by evaluating multiple detectors and integrating deep embedding models.
YOLO was the best detector by mean IoU and produced the strongest crops for pipeline recognition.
Classical methods improved modestly after cropping, while deep methods (ArcFace and FaceNet) achieved near-perfect identification, substantially outperforming handcrafted features.

% No bibliography required unless you add citations.
% \bibliographystyle{IEEEtran}
% \bibliography{report}

\end{document}
